% =====================================================================
%  Multi-Service Ecosystem Collapse Assessment of the Andaman & Nicobar
%  Islands using Remote-Sensing Data Fusion (2001-2024)
%
%  Compile with:  pdflatex main.tex (twice for cross-references)
% =====================================================================
\documentclass[12pt,a4paper]{article}

% ---- AMS math first, then Times-like font family --------------------
% Order matters: newtxmath redefines \Bbbk and a few other blackboard
% symbols, so amssymb must be loaded BEFORE newtxmath (or omitted —
% newtxmath already provides every symbol amssymb does). Here we load
% amsmath + amsfonts only, and let newtxmath supply the AMS symbols.
\usepackage{amsmath,amsfonts}
\usepackage{newtxtext,newtxmath}

% ---- Geometry, spacing, language ------------------------------------
\usepackage[a4paper,margin=1in]{geometry}
\usepackage{setspace}
\onehalfspacing
\usepackage[english]{babel}
\usepackage{microtype}
\usepackage{parskip}

% ---- Graphics, tables, math -----------------------------------------
\usepackage{graphicx}
\usepackage{float}
\usepackage{booktabs}
\usepackage{longtable}
\usepackage{array}
\usepackage{multirow}
\usepackage{caption}
\usepackage{siunitx}
\sisetup{group-separator={,},group-minimum-digits=4}

% ---- Colour palette: an environmental-sciences theme ----------------
% Forest, leaf, sea, earth and accent reds chosen to fit a tropical
% ecosystem-services report. Used for section headings, citation
% links, caption labels and table-header rows.
\usepackage{xcolor}
\definecolor{forestgreen}{HTML}{1B5E20}
\definecolor{leafgreen}  {HTML}{2E7D32}
\definecolor{seablue}    {HTML}{0D47A1}
\definecolor{skyblue}    {HTML}{1976D2}
\definecolor{earthbrown} {HTML}{5D4037}
\definecolor{accentred}  {HTML}{B71C1C}
\definecolor{sand}       {HTML}{F1EAD7}
\definecolor{lightleaf}  {HTML}{E8F1E4}

% Caption labels in forest green and bold, plus a thin colour rule.
\captionsetup{font=small,labelfont={bf,color=forestgreen},
              labelsep=period,textfont=small}

% ---- References (numbered, square brackets) -------------------------
\usepackage[numbers,square,sort&compress]{natbib}

% ---- Hyperlinks: coloured for an environmental-sciences look --------
%   citecolor: leaf green (in-text [n] citations)
%   linkcolor: forest green (TOC, section refs, fig/table refs)
%   urlcolor : sea blue (external URLs and DOI links)
\usepackage[colorlinks=true,
            linkcolor=forestgreen,
            citecolor=leafgreen,
            urlcolor=seablue]{hyperref}

% ---- Section header spacing & colour --------------------------------
\usepackage{titlesec}
\titleformat{\section}
            {\Large\bfseries\color{forestgreen}}
            {\thesection}{1em}{}
\titleformat{\subsection}
            {\large\bfseries\color{leafgreen}}
            {\thesubsection}{1em}{}
\titleformat{\subsubsection}
            {\normalsize\bfseries\color{earthbrown}}
            {\thesubsubsection}{1em}{}

\graphicspath{{figures/}}

% ---- Table-row colouring (booktabs-compatible) ----------------------
\usepackage{colortbl}
\arrayrulecolor{forestgreen!60}

% =====================================================================
\begin{document}

% ---- TITLE PAGE -----------------------------------------------------
\begin{titlepage}
\centering
\vspace*{1.0cm}
{\color{forestgreen}\rule{\linewidth}{1.5pt}}\\[0.2cm]
{\color{leafgreen}\rule{\linewidth}{0.5pt}}\\[0.8cm]
{\Large\bfseries\color{forestgreen}
Quantifying Changes in Ecosystem Services in the\\
Andaman \& Nicobar Islands due to Forest Conversion\\
using Open Global Carbon and Biomass Datasets\par}
\vspace{0.6cm}
{\color{leafgreen}\rule{\linewidth}{0.5pt}}\\[0.2cm]
{\color{forestgreen}\rule{\linewidth}{1.5pt}}\\[1.0cm]
{\large\itshape An Independent Study Report\par}
\vspace{0.4cm}
{\normalsize Submitted in partial fulfilment of the requirements of the course\par}
\vspace{1.8cm}
{\large\bfseries\color{earthbrown} Submitted by\par}
\vspace{0.3cm}
{\large Nikhil Singh\par}
{\normalsize (damanimalay@gmail.com)\par}
\vspace{0.4cm}
{\large Mihir Chaudhary\par}
\vspace{1.2cm}
{\large\bfseries\color{earthbrown} Under the Guidance of\par}
\vspace{0.3cm}
{\large Prof. Ramachandra Prasad\par}
\vfill
{\normalsize\color{seablue}\today\par}
\vspace{0.5cm}
\end{titlepage}

% ---- ABSTRACT -------------------------------------------------------
\section*{Abstract}
This report quantifies the multifunctional collapse of ecosystem services
across the Andaman \& Nicobar Islands (ANI) between 2001 and 2024 by
fusing seven independently produced earth-observation products: ESA
WorldCover 2021 land cover, the Global Forest Watch (GFW) annual
forest-loss raster, the GEDI Level-4B aboveground biomass density grid,
the Saatchi pan-tropical biomass map (independent cross-validation
reference), CHIRPS annual precipitation, SoilGrids clay/silt fractions
and the SRTM 30 m digital elevation model. Three biophysical ecosystem
services---carbon storage, soil retention and habitat quality---are
modelled at 30 m using a combination of GFW $\times$ GEDI carbon
accounting, the Revised Universal Soil Loss Equation (RUSLE) and an
InVEST-style threat-weighted Habitat Quality Index. The three service
losses are normalised, summed and thresholded to produce an Ecosystem
Collapse Index (ECI) identifying pixels where every service is
simultaneously degraded. Over 2001--2024 the islands lost approximately
\SI{19502}{\hectare} of forest, releasing \SI{1390}{\giga\gram} of
CO\textsubscript{2}-equivalent (95\,\% CI 936--1887), with a
statistically significant declining annual trend (Mann--Kendall
$Z=-2.51$, $p=0.012$; Sen's slope $-2.67$ Gg\,CO$_2$e\,yr$^{-1}$). RUSLE
soil-loss rates vary by three orders of magnitude across land covers,
from a median of \SI{2.8}{\tonne\per\hectare\per\year} under closed
canopy to \SI{425}{\tonne\per\hectare\per\year} on bare/sparse surfaces.
Spatial autocorrelation (Moran's $I=0.33$, $p=0.005$ on habitat-quality
deltas; $I=0.20$, $p=0.005$ on RUSLE deltas) confirms that degradation
is spatially clustered rather than random. A supplementary module
quantifies four further services---coastal protection by mangroves,
freshwater yield, pollination suitability and soil organic carbon (SOC)
loss---raising the projected business-as-usual economic damage for 2040
from a lower-bound \$502\,million to an upper-bound \$594\,million. The
combined evidence indicates a regionally contained but spatially
clustered collapse signal that warrants targeted intervention rather
than uniform island-wide management.

\bigskip\noindent\textbf{Keywords:} Ecosystem services, RUSLE, habitat
quality, Ecosystem Collapse Index, GEDI L4B, Andaman \& Nicobar Islands,
Mann--Kendall, Sen's slope, Moran's I, remote sensing.

\clearpage
\tableofcontents
\clearpage

% ---- List of Figures and Tables -------------------------------------
% 25 figures + 4 tables — a navigational index up front is helpful for
% a long report. Both lists adopt the same forest-green colouring as
% the rest of the section headings.
\renewcommand{\listfigurename}{List of Figures}
\renewcommand{\listtablename}{List of Tables}
\listoffigures
\addcontentsline{toc}{section}{List of Figures}
\bigskip
\listoftables
\addcontentsline{toc}{section}{List of Tables}
\clearpage

% =====================================================================
% 1. INTRODUCTION
% =====================================================================
\section{Introduction}

Tropical island ecosystems concentrate disproportionate biological,
hydrological and biogeochemical value into very small areas, and their
exposure to coupled drivers of change---climate variability, land-use
conversion, coastal storms, sediment redistribution---makes them among
the most informative natural laboratories for studying multifunctional
ecosystem degradation. The Andaman \& Nicobar archipelago, comprising
roughly \SI{8249}{\square\kilo\metre} of land scattered across the
eastern Bay of Bengal, contains some of the most intact closed-canopy
tropical forest in South Asia together with a globally significant
mangrove belt, yet the islands have nevertheless experienced
measurable forest loss, coastal disturbance and infrastructure
expansion through the two decades for which a coherent
satellite-observation record exists. Quantifying how much ecosystem
service has actually been lost, where the losses are most concentrated,
whether they aggregate into multifunctional collapse and what the
economic implications might be is therefore both scientifically
tractable and policy-relevant.

\subsection{Literature Survey}

The classical framework for quantifying soil erosion at landscape
scale is the Revised Universal Soil Loss Equation \citep{renard1997}
that decomposes annual soil loss into the product of a rainfall
erosivity factor (R), a soil erodibility factor (K), a slope-length
and steepness factor (LS), a cover-management factor (C) and a
support-practice factor (P). Although the equation was originally
calibrated on rectangular plots in the United States, subsequent
pan-tropical and pan-European applications, notably the high-resolution
RUSLE-based products of \citet{panagos2015} and the global
extrapolation of \citet{borrelli2020}, have demonstrated that the
empirical structure transfers acceptably to remote-sensing inputs when
the factor lookup tables are re-parameterised for tropical conditions.
The C-factor lookup adopted in this study draws on the same
\citet{borrelli2020} tropical defaults, with mangrove and herbaceous
wetland values supplemented from \citet{wischmeier1978}.

Aboveground biomass mapping has shifted decisively in the past five
years from synthetic-aperture-radar-only or optical-only products to
spaceborne lidar. The Global Ecosystem Dynamics Investigation (GEDI)
mission \citep{dubayah2020} provides full-waveform observations on the
International Space Station between $\pm 51.6^\circ$ latitude, and its
Level-4B (L4B) gridded product reports mean aboveground biomass
density (AGBD) at 1 km using random-forest models calibrated against
field plots \citep{dubayah2022}. Cross-validation against the older,
GLAS-calibrated 1 km map of \citet{saatchi2011} reveals systematic
disagreement in dense closed-canopy regions of Africa, South-East Asia
and tropical Asia, with GEDI L4B consistently estimating higher AGBD;
this is documented across several inter-comparison studies and is
attributed both to vintage differences (2003--2007 versus 2019--2023)
and to the conservative nature of the Saatchi map in tall canopies.

Annual deforestation has been operationally tracked since the
publication of the global, 30 m Hansen forest-change product
\citep{hansen2013}, which combines Landsat-derived tree-cover-loss
and annual loss-year layers and is updated by World Resources
Institute via the Global Forest Watch platform. Carbon accounting
applications routinely multiply Hansen lossyear pixels by a baseline
biomass density to estimate annual emissions, with several
international synthesis papers \citep{harris2021} demonstrating that
the approach captures the broad national patterns observed by
independent inventories. The trend-detection literature relies almost
universally on the non-parametric Mann--Kendall test
\citep{mann1945, kendall1975} together with the Theil--Sen slope
estimator \citep{theil1950, sen1968} because annual deforestation is
episodic, non-normal and heavily influenced by individual cyclone or
fire years.

Habitat quality at landscape scale is most commonly modelled with the
InVEST framework \citep{sharp2020}, which assigns each land-cover
class an intrinsic habitat suitability between 0 and 1 and degrades
it as a function of distance to user-specified threats (roads,
settlements, agricultural edges). InVEST's underlying mathematics
follows the half-saturation degradation model of \citet{nelson2009},
and its outputs have been benchmarked against species-distribution
models in tropical settings by \citet{terrado2016}. In this report a
simplified InVEST-style index is computed directly on the ESA
WorldCover 2021 land-cover map without an external threat layer, on
the grounds that the ANI archipelago has relatively few high-impact
linear infrastructure features outside the south Andaman urban
cluster.

Multi-service synthesis indices that aggregate independently modelled
service losses into a single decision-relevant raster have a longer
history in the freshwater literature, but the conceptual approach used
here---normalise each service to $[0,1]$, sum the normalised harms
and threshold at the 95\textsuperscript{th} percentile to identify
``triple-collapse'' pixels---is consistent with the multi-criteria
hotspot synthesis used by \citet{egoh2008} and the bundle/anti-bundle
analysis of \citet{rao2018}. Spatial autocorrelation in the resulting
hotspot raster is quantified using the global Moran's $I$ statistic
\citep{moran1950} with a queen-case 3$\times$3 neighbourhood and a
permutation-based test of significance, following the spatial-
econometrics conventions established for ecosystem-service mapping in
\citet{anselin1995}.

Cross-validation of remote-sensing biomass products requires more than
a Pearson correlation, because two products can rank pixels similarly
yet sit on completely different absolute scales. \citet{lin1989}
introduced the Concordance Correlation Coefficient (CCC) that
multiplies Pearson's $r$ by a bias-correction factor measuring the
displacement of the regression line from the 1:1 diagonal; products
that disagree systematically yield very low CCC even when $r$ is
moderate. Reduced Major Axis (RMA) regression \citep{warton2006}
complements ordinary least squares by treating both variables as
having error, which is the realistic assumption when validating one
remote-sensing model against another rather than against ground
truth. Bootstrap confidence intervals on RMSE, bias, $r$ and CCC are
standard practice in modern biomass cross-validation
\citep{efron1993}.

The economic translation of biophysical losses into a single
willingness-to-pay or social-cost figure typically applies the United
States Environmental Protection Agency social cost of CO$_2$
\citep{epa2023} multiplied by the CO$_2$-equivalent emissions
estimate, with the well-known caveat that any single SCC figure
captures only one decision-maker's value function. For coastal
protection, the global willingness-to-pay calculations of
\citet{menendez2020} provide a defensible per-hectare per-year
mangrove buffering value, while soil-organic-carbon losses from
forest-to-cropland conversion are summarised by the meta-analytic
factor of \citet{don2011}.

\subsection{Objectives of the Study}

The objectives of this study are: (i) to assemble a coherent
multi-service remote-sensing dataset for the Andaman \& Nicobar
archipelago at 30 m spatial resolution and to confirm that the
constituent rasters are mutually projectable, mutually clippable and
mutually comparable through the 2001--2024 reference period; (ii) to
quantify annual carbon-equivalent losses from forest deforestation by
spatially intersecting the GFW annual forest-loss layer with the GEDI
L4B aboveground biomass density grid and to test the resulting time
series for monotonic trends using the Mann--Kendall and Theil--Sen
statistics; (iii) to compute pixel-level RUSLE soil-loss rates and
their land-cover-class statistics, including the counterfactual delta
relative to a forested baseline; (iv) to estimate an InVEST-style
habitat quality index per pixel and to summarise its distribution per
land-cover class, including the change relative to a pre-disturbance
baseline; (v) to synthesise the three service losses into a single
Ecosystem Collapse Index, identify triple-collapse hotspots at the
95\textsuperscript{th} percentile, and test the resulting hotspot
pattern for spatial autocorrelation using a permutation-based Moran's
$I$; (vi) to cross-validate the GEDI L4B biomass values against the
independent Saatchi 2011 pan-tropical product on the forest-only
pixel set using Pearson's $r$, the coefficient of determination,
Lin's CCC, RMSE, bias, OLS regression and RMA regression with
bootstrap-derived confidence intervals; (vii) to extend the analysis
with a supplementary module covering coastal protection by
mangroves, Budyko-style freshwater yield, Lonsdorf-style pollination
suitability and SoilGrids-derived organic-carbon loss; and (viii) to
translate the combined biophysical losses into projected economic
damages to 2040 under business-as-usual, conservation and escalation
scenarios.

\subsection{Study Area}

The Andaman \& Nicobar Islands form an arcuate Union Territory of
India scattered between approximately $6.5^\circ$N and $13.5^\circ$N
in the eastern Bay of Bengal, with the Andaman group in the north
and the Nicobar group separated from it by the ten-degree channel.
Total land area is approximately \SI{8249}{\square\kilo\metre},
of which closed tropical broadleaf forest accounts for roughly
\SI{6390}{\square\kilo\metre} (Tree Cover ESA class 10) and mangrove
forest a further \SI{590}{\square\kilo\metre} (ESA class 95). The
remaining surface comprises grassland, scattered cropland, built-up
settlements concentrated around Port Blair, and bare or sparse
substrate primarily on the more arid Nicobar interior. Annual mean
precipitation in the CHIRPS record is \SI{3400}{\milli\metre} on Tree
Cover and exceeds \SI{3840}{\milli\metre} on the Mangrove fringe, and
the SRTM 30 m DEM shows that the majority of the archipelago lies
below \SI{100}{\metre} elevation, with the dominant relief being the
Saddle Peak ridge in north Andaman and a smaller central ridge on
Great Nicobar. The islands have experienced major natural
disturbances during the study period---most notably the December 2004
Indian Ocean earthquake and tsunami---and these disturbances are
reflected in the GFW forest-loss series and the carbon-loss spikes
documented in Section~\ref{sec:carbon}.

% =====================================================================
% 2. MATERIALS AND METHODS
% =====================================================================
\section{Materials and Methods}

\subsection{Data Used}
\label{sec:data}

Seven primary remote-sensing rasters were used in this study, each
clipped to the ANI administrative boundary, reprojected to UTM Zone
46N (EPSG:32646) and resampled to a common 30 m pixel grid where
necessary. Table~\ref{tab:datasets} summarises the data sources,
their native resolutions and the post-processing applied. The
processed copies are stored in \texttt{data/processed/} and are
read by every analysis script.

\begin{table}[H]
\centering
\caption{Earth-observation and ancillary datasets used in this study.}
\label{tab:datasets}
\small
\begin{tabular}{p{4.0cm} p{3.3cm} p{1.8cm} p{5.5cm}}
\toprule
\rowcolor{forestgreen!15}
\textbf{Dataset} & \textbf{Source} & \textbf{Native res.} & \textbf{Use}\\
\midrule
ESA WorldCover 2021 & ESA Climate Change Initiative & 10 m & Land-cover stratification, RUSLE C-factor lookup, InVEST habitat suitability, masking. \\
GFW Forest Loss 2001--2023 & Hansen / WRI Global Forest Watch & 30 m & Annual deforestation events, drives carbon-loss accounting. \\
GFW Tree Cover 2000 baseline & Hansen / WRI GFW & 30 m & Pre-2001 canopy reference. \\
GEDI L4B AGBD & NASA GEDI mission v2.1 & 1 km (gridded) & Aboveground biomass density baseline; carbon-stock input. \\
Saatchi 2011 Pan-Tropical AGB & \citet{saatchi2011} & 1 km & Independent cross-validation reference. \\
CHIRPS annual precipitation & UCSB Climate Hazards Center & $\sim$5 km & RUSLE R-factor (rainfall erosivity), freshwater-yield module. \\
SoilGrids Clay+Silt fraction & ISRIC SoilGrids v2 (via GEE) & 250 m & RUSLE K-factor proxy and SOC baseline. \\
SRTM 30 m DEM & USGS / NASA SRTM v3 & 30 m & Slope-length and steepness (LS) factor, coastal-elevation factor. \\
ANI administrative boundary & Survey of India shapefile & vector & Region-of-interest clipping. \\
\bottomrule
\end{tabular}
\end{table}

All rasters were resampled to the ESA WorldCover 30 m grid using
nearest-neighbour interpolation for the categorical layers (ESA,
GFW lossyear) and bilinear interpolation for the continuous layers
(GEDI, Saatchi, CHIRPS, SoilGrids, SRTM). Nodata cells were
propagated through every analytical step using NumPy NaN handling so
that downstream statistics never mistake ocean for unmodelled land.

\subsection{Carbon-Loss Accounting from GFW $\times$ GEDI}
\label{sec:methods-carbon}

Annual carbon-equivalent losses were computed by intersecting the
GFW \emph{lossyear} raster $L(x,y)\in\{0,1,\dots,23\}$ with the GEDI
L4B biomass density grid $B(x,y)$ in Mg\,ha$^{-1}$. For each year
$t$ from 2001 to 2023 a binary mask $M_t = \mathbb{1}\{L = t-2000\}$
was extracted and the lost aboveground biomass in that year was
computed pixel-by-pixel as

\[
B_t(x,y) \;=\; M_t(x,y) \,\cdot\, B(x,y) \,\cdot\, 0.09
\]

with the constant 0.09 converting per-hectare biomass density to
per-pixel biomass for 30 m pixels. Aboveground carbon $C_t$ was
obtained by multiplying $B_t$ by 0.47, the standard biomass-to-carbon
fraction \citep{westlake1963}, and converted to CO$_2$ equivalent by
multiplying by $44/12$:

\[
\text{CO}_2\text{e}_t \;=\; 0.47 \cdot \frac{44}{12} \cdot B_t
\;=\; 1.723 \cdot B_t .
\]

Aggregating over $(x,y)$ gives the annual time series shown in
Figure~\ref{fig:trend}. The cumulative gain CO$_2$e from the GFW
\emph{gain} layer is computed in the same way using a forested-pixel
assumption of \SI{96.7}{\mega\gram\of\dryMatter\per\hectare} for
recovering canopy, and the net balance time series is included in
the carbon-results CSV. The 24-year total CO$_2$e loss together with
its 95\,\% confidence interval is reported in
Table~\ref{tab:headline}.

\subsection{Trend Detection: Mann--Kendall and Theil--Sen}

The annual CO$_2$e and annual lost-area time series were tested for
monotonic trends using the non-parametric Mann--Kendall (MK) test
\citep{mann1945, kendall1975} with the tie-corrected variance
$\sigma_S^2 = [n(n-1)(2n+5) - \sum_j t_j(t_j-1)(2t_j+5)] / 18$ and a
continuity correction giving

\[
Z = \begin{cases}
(S-1)/\sigma_S, & S > 0 \\
0, & S = 0 \\
(S+1)/\sigma_S, & S < 0 .
\end{cases}
\]

The Theil--Sen median slope was computed as the median of all
pairwise slopes $(y_j - y_i)/(t_j - t_i)$ for $i<j$, with the
95\,\% confidence interval derived from the Hollander--Wolfe
normal approximation $c_\alpha = 1.96\sqrt{\sigma_S^2}$ and the
sorted-slope positions $\lfloor (N_{\text{pairs}}-c_\alpha)/2\rfloor$
and $\lceil (N_{\text{pairs}}+c_\alpha)/2\rceil$. The Theil--Sen line
is anchored at the data centroid $(\widetilde{t},\widetilde{y})$ for
plotting (Figure~\ref{fig:trend}); the slope-uncertainty ribbon
shown in the figure visualises slope-CI lines pivoting through this
centroid and is therefore a slope envelope rather than a predictive
interval.

\subsection{RUSLE Soil-Loss Modelling}
\label{sec:methods-rusle}

Annual soil loss $A$ in t\,ha$^{-1}$\,yr$^{-1}$ was computed as
$A = R\cdot K\cdot LS\cdot C\cdot P$. The R-factor was derived from
CHIRPS annual precipitation using the regional regression $R = 81.5 +
0.38\,P_{annual}$ \citep{panagos2015}. The K-factor was approximated
from the SoilGrids clay+silt fraction using the simplified relation
$K = 0.005 + 0.0006(\text{clay}+\text{silt})$, clipped to a plausible
physical range of $[0.005, 0.060]$. The LS-factor was computed from
the SRTM DEM using the \citet{moore1992} flow-accumulation
formulation with a steep-slope cap of 50 to avoid cliff-edge
artefacts. The C-factor was a per-class lookup
(Table~\ref{tab:cfactor}) drawing on the tropical defaults of
\citet{borrelli2020}. The P-factor was set uniformly to 1.0 (no
active management practice information available at 30 m). Pixel
losses were aggregated by ESA WorldCover class to produce
Table~\ref{tab:rusle-by-class} and the violin distribution in
Figure~\ref{fig:rusle-violin}. The four component factors are
visualised over both island groups in
Figure~\ref{fig:rusle-factors}, where the K-factor is rendered on the
high-contrast \texttt{cividis} colormap with 2--98\,\% percentile
stretch and the C-factor on \texttt{magma\_r} with a log-norm so the
98\,\% tree-cover background does not visually swamp the rare
high-erosion cropland and bare patches.

\begin{table}[H]
\centering
\caption{RUSLE C-factor lookup applied per ESA WorldCover class
\citep{borrelli2020,wischmeier1978}.}
\label{tab:cfactor}
\small
\begin{tabular}{l l c}
\toprule
\rowcolor{forestgreen!15}
\textbf{Class code} & \textbf{Land cover} & \textbf{C-factor} \\
\midrule
10 & Tree cover     & 0.001 \\
20 & Shrubland      & 0.035 \\
30 & Grassland      & 0.060 \\
40 & Cropland       & 0.280 \\
50 & Built-up       & 0.000 \\
60 & Bare / sparse  & 0.450 \\
90 & Wetland (herb.)& 0.020 \\
95 & Mangroves      & 0.010 \\
\bottomrule
\end{tabular}
\end{table}

\subsection{InVEST-Style Habitat Quality Index}
\label{sec:methods-habitat}

Each ESA WorldCover class was assigned an intrinsic habitat
suitability score $H_c \in [0,1]$ drawing on the InVEST tropical
defaults of \citet{sharp2020}: Tree cover 1.00, Mangroves 0.90,
Shrubland 0.50, Wetland 0.45, Grassland 0.38, Bare 0.15, Cropland
0.10, Built-up 0.00. The score was applied per pixel and modulated by
a degradation function $D(x,y)$ that decays linearly with distance
from threat sources (settlements, infrastructure edges) where such
features were derivable from the land-cover map itself. The final
habitat quality $Q(x,y)$ was

\[
Q(x,y) = H_c(x,y) \cdot \left[1 - \frac{D(x,y)^k}{D(x,y)^k + D_{1/2}^k}\right]
\]

with the half-saturation constant $D_{1/2}$ set to 0.5 and the shape
parameter $k=2.5$, following the original InVEST default. The change
$\Delta Q$ relative to a pre-disturbance counterfactual---in which
all currently disturbed pixels are reassigned to their pre-loss tree
cover or mangrove state---is the diagnostic used in the synthesis.
Per-class distributions are visualised in
Figure~\ref{fig:hq-violin}.

\subsection{Ecosystem Collapse Index (ECI)}
\label{sec:methods-eci}

Three normalised harm layers were combined into a single ECI raster.
The habitat harm layer $h_{\text{hab}}$ was $-\Delta Q$ clipped to
$[0,1]$, the soil harm layer $h_{\text{soil}}$ was the log-transformed
RUSLE delta normalised by its 99.9\textsuperscript{th} percentile,
and the carbon harm layer $h_{\text{C}}$ was a binary flag indicating
any GFW forest-loss event during 2001--2023. The combined index was

\[
\text{ECI}(x,y) = \frac{h_{\text{hab}}(x,y) + h_{\text{soil}}(x,y)
                       + h_{\text{C}}(x,y)}{3}
\]

producing values in $[0,1]$. Pixels exceeding the
95\textsuperscript{th} percentile of the impacted-pixel distribution
were flagged as triple-collapse hotspots (Figure~\ref{fig:eci-map};
Figure~\ref{fig:hexbin}). Global Moran's $I$ was computed on the
binary hotspot mask with a queen-case 3$\times$3 neighbourhood
($w_{ij}=1/8$ for adjacent cells, zero elsewhere) using the
convolution formulation, with a 199-replicate permutation null
distribution providing the empirical $p$-value.

\subsection{Cross-Validation: GEDI L4B versus Saatchi 2011}
\label{sec:methods-cv}

The independent Saatchi pan-tropical AGB map was used to cross-
validate the GEDI L4B aboveground biomass values pixel-by-pixel after
restricting both rasters to forest-only land-cover classes
(ESA 10 Tree cover and ESA 95 Mangroves). Restriction is essential
because Saatchi is a forest-only product whose nominal values over
non-forest pixels are uninformative, while GEDI L4B's gap-filled grid
returns numerical estimates everywhere. With \num{2314513} forest
pixels in the overlap, six agreement metrics were computed: Pearson's
$r$, the coefficient of determination $R^2$, Lin's Concordance
Correlation Coefficient $\rho_c$ \citep{lin1989}, RMSE, bias
(GEDI~$-$~Saatchi) and mean absolute error. Two regression lines were
fitted: ordinary least squares with Saatchi as predictor, and Reduced
Major Axis regression $\hat{b}_{RMA} = \mathrm{sign}(r)\,\sigma_y/
\sigma_x$ \citep{warton2006} which is the appropriate fit when both
variables carry error. Bootstrap 95\,\% confidence intervals on every
metric were derived from 500 resamples of the forest-pixel subsample.

\subsection{Supplementary Ecosystem-Service Modules}
\label{sec:methods-supp}

Four lightweight supplementary services were computed to bracket
the lower-bound carbon-plus-soil-plus-habitat damage figure. Coastal
protection was modelled as a mangrove buffer index $CPI =
\mathbb{1}_{\text{mangrove}} \cdot f_{\text{elev}} \cdot f_{\text{coast}}$
where $f_{\text{elev}}$ decays linearly from 1 below \SI{10}{\metre}
elevation to 0 at \SI{30}{\metre} and $f_{\text{coast}}$ decays
linearly from 1 within \SI{1}{\kilo\metre} of water to 0 at
\SI{5}{\kilo\metre}, following \citet{menendez2020}. The per-hectare
value was set at \$2240/ha/yr (10\,\% of the global Menendez et al.
willingness-to-pay benchmark, retained as a conservative ANI-scaled
lower bound), and the 15-year net present value of high-value
($CPI > 0.5$) area lost was computed using the annuity factor
$a = (1 - (1+r)^{-n})/r$ at $r=0.04$, $n=15$ giving $a \approx 11.12$
(corrected from a previous undiscounted version). Freshwater yield
was estimated as $Y = P\cdot(1-\kappa_c)$ where $\kappa_c$ is a
class-specific evapotranspiration coefficient drawn from
\citet{zhang2001,sun2006}. Pollination suitability followed the
Lonsdorf-style index $S = \sqrt{F\cdot N}$ with class-specific
floral $F$ and nesting $N$ scores from the tropical defaults of
\citet{koh2016}. Baseline soil organic carbon was modelled hybridly:
where the SoilGrids clay+silt raster has valid data, a pedotransfer
formula $\text{SOC} = 40 + 0.5\,(\text{clay}+\text{silt})$ was
applied; where SoilGrids reports zero (a coarse-to-fine resampling
artefact), a literature class-mean SOC was used following
\citet{powers2018,donato2011}. SOC loss from deforested pixels was
computed using the Don et al. \citep{don2011} 25\,\% loss factor and
converted to CO$_2$e and US dollar damage at the US EPA 2023 social
cost of carbon of \$51/tCO$_2$.

\subsection{Economic Forecasting to 2040 and 2060}

The 2040 business-as-usual economic damage forecast combined the
historic deforestation velocity (\SI{812.6}{\hectare\per\year}, the
average annual loss across 2001--2023) with three damage categories:
CO$_2$e damage at the social cost of carbon (\$51/tCO$_2$), dredging
cost from sediment loss (\$5/tonne removed) and ecosystem
habitat-replacement cost (\$3000/ha for tropical forest). The
upper-bound version added the coastal-protection NPV and SOC damage
from the supplementary module. Three scenarios were extended to 2060:
a Conservation scenario assuming zero new loss after 2024, a
Business-As-Usual scenario continuing the historic velocity, and an
Escalation scenario doubling the velocity.

% =====================================================================
% 3. RESULTS AND DISCUSSION
% =====================================================================
\section{Results and Discussion}

\subsection{GEDI Aboveground Biomass Baseline}
\label{sec:carbon-baseline}

Before the loss budget could be computed, an aboveground biomass (AGB)
baseline had to be established at every land pixel of the archipelago.
The NASA GEDI Level-4B 1\,km gridded AGB density product served as
this baseline because it is the only spaceborne lidar product that
directly measures the vertical canopy structure responsible for the
bulk of the biomass signal, and because its random-forest model is
calibrated against field plots at the footprint level rather than
relying on optical proxies. The product was bilinearly resampled to
the 30\,m ESA WorldCover grid so that every analytical pixel had a
biomass estimate.

Figure~\ref{fig:agb-baseline} renders the resulting biomass surface
separately for the Andaman and Nicobar groups, overlaid with the GFW
forest-gain pixels (dilated $\times 2$ for visibility). The Andaman
mean AGB is 208 Mg\,ha$^{-1}$ with a 95\textsuperscript{th} percentile
of 336 Mg\,ha$^{-1}$; the Nicobar mean is 184 Mg\,ha$^{-1}$ with a
95\textsuperscript{th} percentile of 322 Mg\,ha$^{-1}$. The
spatial pattern is broadly what one would expect for an undisturbed
tropical archipelago: the densest biomass is on the central spine of
North and Middle Andaman, on the Saddle Peak ridge, and on the deep
interior of Great Nicobar; the lowest values occur on the coastal
fringe (mangrove-dominated, lower per-pixel AGB than upland tree
cover), on the small Nicobar islets where canopy structure is more
fragmented, and around the few inhabited settlements. Forest-gain
patches (magenta dots) are scattered preferentially along the eastern
coast of Middle Andaman---a known regeneration corridor following the
2004 tsunami---and on the southern tip of Great Nicobar.

\begin{figure}[H]
\centering
\includegraphics[width=0.92\textwidth]{carbon/agb_gedi_baseline_map.png}
\caption{GEDI L4B aboveground biomass density (Mg\,ha$^{-1}$) over the
Andaman (left) and Nicobar (right) island groups, rendered on a green
sequential ramp. Forest-gain pixels from the GFW 2000--2012 product,
dilated by a factor of two for visibility, are overlaid in magenta.
Inset boxes show the mean AGB and the 95\textsuperscript{th}
percentile for each group.}
\label{fig:agb-baseline}
\end{figure}

This biomass baseline drives every downstream carbon-loss calculation
in Section~\ref{sec:carbon}, and the absolute scale of the GEDI L4B
estimates is independently cross-validated against the Saatchi 2011
product in Section~\ref{sec:cv-results} below. The reader should bear
in mind that any uncertainty in the GEDI baseline propagates
multiplicatively into the carbon-loss figure: a 10\,\% systematic
overestimate of mean AGB inflates the 24-year CO$_2$e loss by the
same 10\,\%. The bootstrap 95\,\% confidence interval reported below
(936--1887 Gg CO$_2$e) is dominated by the spatial variance of the
biomass field rather than by uncertainty in the temporal
deforestation count.

\subsection{Carbon-Loss Trajectory and Trend Significance}
\label{sec:carbon}

Across the 24-year analysis window the archipelago lost
approximately \SI{19502}{\hectare} of forest, equivalent to
\SI{1390}{\giga\gram} CO$_2$e (95\,\% bootstrap CI 936--1887). The
annual time series is visibly episodic
(Figure~\ref{fig:trend}): three pronounced spikes occur in 2005
($\sim$\SI{180}{\giga\gram}), 2008 ($\sim$\SI{125}{\giga\gram}) and
2020 ($\sim$\SI{160}{\giga\gram}), separated by years of low loss.
The Mann--Kendall test on the annual CO$_2$e series returns
$Z=-2.51$, $p=0.0122$ with the decreasing-trend direction, and
Theil--Sen estimates a slope of
$-2.67$ Gg\,CO$_2$e\,yr$^{-1}$ (95\,\% CI $-5.25$ to $-0.66$). The
parallel test on annual lost area gives an even stronger signal,
$Z=-3.35$, $p=0.0008$, Sen slope $-54.1$ ha\,yr$^{-1}$ (95\,\% CI
$-83.8$ to $-20.4$). The slope confidence interval excludes zero in
both cases, so the statistical declining trend is robust to the
specification.

\begin{figure}[H]
\centering
\includegraphics[width=0.92\textwidth]{validation/carbon_loss_trend_with_ci.png}
\caption{Annual CO$_2$e loss in ANI between 2001 and 2024 (red bars),
Theil--Sen median slope (dashed blue) and the slope-envelope 95\,\%
confidence band (light blue, clipped at zero because CO$_2$e loss is
strictly non-negative). The Mann--Kendall test annotation in the
upper-left reports a statistically significant declining trend
($Z=-2.51$, $p=0.012$). The three large spikes around 2005, 2008 and
2020 dominate the trend's anchor points.}
\label{fig:trend}
\end{figure}

A companion display of the same series with cumulative emissions
overlaid as a green twin-axis curve and broken down by within-year
hectares of forest loss is shown in
Figure~\ref{fig:carbon-timeseries}. The cumulative curve is what
matters for the policy question of \emph{how much carbon has already
left the atmosphere-canopy reservoir}: every single spike adds a
permanent step that the subsequent quieter years do not erase, so by
the end of 2024 the islands sit at a cumulative emission of
$\approx$\SI{1390}{\giga\gram} CO$_2$e regardless of the declining
annual trend in the bar series. This is a crucial framing distinction
because a casual reader of the trend figure could conclude that the
crisis is ``getting better''; the cumulative perspective makes clear
that the standing debt continues to accrue.

\begin{figure}[H]
\centering
\includegraphics[width=0.92\textwidth]{carbon/carbon_annual_loss_timeseries.png}
\caption{Annual forest loss area (left axis, red bars) and the
cumulative CO$_2$e emission curve (right axis, green line) for ANI
2001--2024. Despite the visibly declining annual loss series, the
cumulative curve continues to climb monotonically, ending at
$\approx$ 1390 Gg CO$_2$e of debt against the atmosphere by 2024.}
\label{fig:carbon-timeseries}
\end{figure}

It is important not to over-interpret the declining MK trend, however.
The annual series has only 24 observations, of which three (2005,
2008, 2020) account for nearly 40\,\% of the cumulative loss. The
2005 spike is consistent with documented post-2004-tsunami
reconstruction and incidental clearing pulses, the 2008 spike with a
period of intensified illegal logging investigations reported by the
Forest Department, and the 2020 spike with cyclone-related damage
during the COVID-19 administrative window. In other words, the
``declining trend'' is more accurately framed as ``fewer large-event
years recently'' than as a steady reduction in deforestation
pressure.

\subsection{Where Did the Loss Occur? Spatial Carbon Hotspots}

Beyond the temporal pattern, knowing \emph{where} on the archipelago
carbon emissions occurred matters for both root-cause diagnosis and
intervention design. Figure~\ref{fig:carbon-hotspots} maps the
spatially-resolved CO$_2$e loss intensity by aggregating per-pixel
losses across all study years. Two patterns dominate. First, the
Andaman emissions are concentrated on Middle and South Andaman around
the existing road network and settlement boundaries, with a thin
band of high-intensity loss along the eastern coast of North Andaman
that aligns with documented tsunami-period reconstruction sites.
Second, the Nicobar emissions are dominated by a single contiguous
hotspot on the south-eastern flank of Great Nicobar---a pattern
that pre-dates the headline-grabbing 2020s development proposals and
reflects steady, low-intensity clearing through the 2000s and 2010s.
This spatial concentration is precisely what makes the Moran's $I$
test return a strongly clustered signature (Section~\ref{sec:eci-results}).

\begin{figure}[H]
\centering
\includegraphics[width=0.92\textwidth]{carbon/carbon_loss_hotspots_map.png}
\caption{Spatial aggregation of carbon-loss intensity (cumulative
CO$_2$e per pixel, 2001--2024) for the Andaman (left) and Nicobar
(right) groups. The Andaman pattern follows the road and settlement
network; the Nicobar pattern is dominated by a single contiguous
hotspot on south-east Great Nicobar.}
\label{fig:carbon-hotspots}
\end{figure}

\subsection{Loss versus Gain and the Net Carbon Balance}

The GFW product also reports a separate forest-\emph{gain} layer
covering 2000--2012, which allows the loss totals to be partially
offset by the carbon sequestered in newly regenerating canopy.
Figure~\ref{fig:loss-vs-gain} displays both layers on the same map,
with cumulative loss in red and cumulative gain in yellow. Three
observations follow. (i) Loss and gain are not co-located: the loss
pixels occur preferentially along roads and settlement edges, while
gain pixels are scattered in deeper interior locations not currently
under active conversion pressure. (ii) The gain area at
\SI{5806}{\hectare} (\SI{4256}{\hectare} on the Andamans plus
\SI{1550}{\hectare} on the Nicobars) is roughly one-third of the
loss area at \SI{19502}{\hectare}. (iii) Gain is reported only for
the 2000--2012 window because the GFW gain product has not been
updated for later years, so the net-balance calculation has an
inherent forward-bias: post-2012 regrowth that has matured into
detectable canopy is missing from the gain side of the ledger.

\begin{figure}[H]
\centering
\includegraphics[width=0.92\textwidth]{carbon/carbon_loss_vs_gain_spatial.png}
\caption{Forest loss (red, 2001--2024) versus forest gain (yellow,
2000--2012) for the two island groups, displayed against the
intact-land basemap. Loss exceeds gain by a factor of approximately
3.4. Loss and gain are spatially anticorrelated: clearing occurs on
the road and settlement perimeter while regrowth concentrates in
the deeper interior.}
\label{fig:loss-vs-gain}
\end{figure}

When the gain CO$_2$e is netted against the loss CO$_2$e on an
annual basis, the resulting net-balance time series shown in
Figure~\ref{fig:net-balance} remains firmly negative for every
year of the record. The cumulative net debt at end-2024 is
approximately \SI{882}{\giga\gram} CO$_2$e once the modelled gain
of \SI{508}{\giga\gram} is subtracted from the gross loss of
\SI{1390}{\giga\gram}. The qualitative picture is unchanged: the
archipelago is a net carbon source even after granting full credit
to the GFW gain product.

\begin{figure}[H]
\centering
\includegraphics[width=0.85\textwidth]{carbon/carbon_net_balance_figure.png}
\caption{Annual net carbon balance for ANI 2001--2024 (loss minus
gain, expressed as Gg CO$_2$e). The series is uniformly negative
(loss-dominated) and the cumulative debt at end-2024 is
$\approx$ 882 Gg CO$_2$e after granting full credit to the GFW
forest-gain layer.}
\label{fig:net-balance}
\end{figure}

\subsection{RUSLE Component Maps and Soil-Loss Distributions}
\label{sec:rusle-results}

The four RUSLE factor component maps in Figure~\ref{fig:rusle-factors}
show that rainfall erosivity (R) is highest on the wetter northern
Andamans (CHIRPS annual mean over \SI{3500}{\milli\metre}); soil
erodibility (K) varies modestly along the central Andaman ridge where
clay+silt fractions exceed 50\,\%; the slope-length-and-steepness
factor (LS) is concentrated on the steeper Saddle Peak block and on
the Great Nicobar interior; and the cover-management factor (C),
plotted on a log scale, is dominated by the very low forest baseline
($C=0.001$) with rare high-erosion cropland and bare patches glowing
bright on the otherwise dark background.

\begin{figure}[p]   % own page — figure is 4 rows tall
\centering
\includegraphics[height=0.88\textheight,keepaspectratio]{soil/rusle_factor_components_map.png}
\caption{RUSLE four-factor component maps for the Andaman (left) and
Nicobar (right) island groups. R is rainfall erosivity, K is soil
erodibility, LS is the slope-length-and-steepness factor and C is the
cover-management factor (plotted on a log scale so the rare
high-erosion patches are visible against the 98\,\% tree-cover
background).}
\label{fig:rusle-factors}
\end{figure}

The resulting per-class soil-loss distributions span three orders of
magnitude (Figure~\ref{fig:rusle-violin}, Table~\ref{tab:rusle-by-class}).
Closed-canopy tree cover yields a median loss of just
\SI{2.83}{\tonne\per\hectare\per\year}---well within the soil-formation
rate range of \SI{1}{\tonne\per\hectare\per\year}--\SI{11}{\tonne\per\hectare\per\year} typically cited for
tropical soils \citep{montgomery2007}---while bare/sparse surfaces
exceed \SI{424}{\tonne\per\hectare\per\year} at the median and saturate
the RUSLE physical ceiling of 500 at the upper tail. Mangroves
($\sim$\SI{12.4}{\tonne\per\hectare\per\year}), wetlands
($\sim$\SI{10.5}{\tonne\per\hectare\per\year}) and shrubland
($\sim$\SI{33}{\tonne\per\hectare\per\year}) sit comfortably between
forest and grassland/cropland. Built-up surfaces register zero soil
loss because RUSLE's C-factor collapses for impervious surfaces---a
modelling convention rather than a measurement.

\begin{figure}[H]
\centering
\includegraphics[width=0.78\textwidth]{synthesis/stat_distribution_soil_erosion.png}
\caption{Per-land-cover-class distributions of RUSLE annual soil loss
on a log axis, sorted ascending by class median. Built-up was dropped
because all values are exactly zero under RUSLE's impervious-surface
convention; Shrubland is rendered with a dashed outline because its
sample size ($n=99$) is too small for the violin shape to be
statistically reliable.}
\label{fig:rusle-violin}
\end{figure}

The aggregate RUSLE soil-loss surface itself, computed as the
pixel-by-pixel product of all five factors and capped at the physical
ceiling of 500 t\,ha$^{-1}$\,yr$^{-1}$, is shown in
Figure~\ref{fig:rusle-map}. The mean island-wide soil loss is
\SI{7.9}{\tonne\per\hectare\per\year}---comparable to the
\SI{5}{\tonne\per\hectare\per\year}--\SI{12}{\tonne\per\hectare\per\year}
range reported by \citet{borrelli2020} for closed-canopy tropical
sub-regions---and the maximum reaches the 500-cap on a handful of bare
and cropland pixels. The spatial pattern is closely tied to
land-cover: the islands appear uniformly low-yellow over the dominant
tree-cover background, with discrete high-erosion spots on the
cropland fringe around Port Blair and the bare patches on Great
Nicobar's southern slopes.

\begin{figure}[H]
\centering
\includegraphics[width=0.78\textwidth]{soil/rusle_soil_loss_map.png}
\caption{Pixel-level RUSLE annual soil erosion risk $A = R\cdot K\cdot
LS\cdot C\cdot P$ for the Andaman (left) and Nicobar (right) island
groups. The colour bar is log-scaled to reveal the three-orders-of-
magnitude range between the protected forest background and the
high-erosion cropland and bare patches. Island-wide mean is
\SI{7.9}{\tonne\per\hectare\per\year}.}
\label{fig:rusle-map}
\end{figure}

The \emph{counterfactual delta} between current soil loss and the
all-forested baseline---the soil loss \emph{caused by} land-cover
conversion---is what truly matters for the ecosystem-service
accounting, and is mapped in Figure~\ref{fig:rusle-delta}. The
spatial pattern of the delta is much sharper than the raw soil-loss
map because the forest background contributes effectively zero to the
delta. Hotspots in the delta map align almost perfectly with the
land-cover-conversion footprint: the south Andaman cropland belt
around Port Blair, the small bare patches on Great Nicobar's
southern slopes, and scattered grassland edges throughout the
archipelago. The total ``new'' erosion attributable to forest
conversion sums to roughly \SI{163000}{\tonne\per\year}, of which
74\,\% is contributed by grassland and cropland alone.

\begin{figure}[H]
\centering
\includegraphics[width=0.78\textwidth]{soil/rusle_soil_loss_delta_map.png}
\caption{Counterfactual delta between current RUSLE soil loss and the
all-forested baseline. Hotspots reveal the spatial footprint of
land-cover conversion that is currently increasing sediment yield
relative to a fully forested landscape.}
\label{fig:rusle-delta}
\end{figure}

A complementary horizontal bar chart in
Figure~\ref{fig:rusle-bar} summarises the per-class totals from
Table~\ref{tab:rusle-by-class} in a visual format: tree cover
contributes the largest absolute total in t\,yr$^{-1}$ because its
area is enormous (over half a million hectares), but the per-hectare
\emph{rate} is dominated by cropland and bare/sparse classes that
occupy a tiny fraction of the island area. This is the standard
ecosystem-services trade-off: a small area of intensively-managed or
unvegetated land contributes disproportionate damage on a per-pixel
basis.

\begin{figure}[H]
\centering
\includegraphics[width=0.82\textwidth]{soil/rusle_erosion_by_landcover.png}
\caption{Aggregate RUSLE soil-loss totals per ESA WorldCover class,
showing both the total annual sediment yield (left bars) and the
per-hectare rate (right bars). Tree cover dominates the total
because its area is largest; cropland and bare/sparse dominate the
per-hectare rate.}
\label{fig:rusle-bar}
\end{figure}

\begin{table}[H]
\centering
\caption{RUSLE soil-loss statistics by ESA WorldCover class.}
\label{tab:rusle-by-class}
\small
\begin{tabular}{l r r r r r}
\toprule
\rowcolor{forestgreen!15}
\textbf{Land cover} & \textbf{Area (ha)} & \textbf{Mean} & \textbf{Median} & \textbf{Max} & \textbf{Total} \\
\rowcolor{forestgreen!15}
                    &                   & (t/ha/yr)    & (t/ha/yr)       & (t/ha/yr)    & (t/yr)        \\
\midrule
Tree cover (Forest) & \num{524713} & 3.67   & 2.83   & 41.0  & \num{1927237} \\
Mangroves           & \num{49239}  & 17.34  & 12.41  & 336.4 & \num{853674}  \\
Grassland           & \num{13222}  & 87.12  & 56.36  & 500.0 & \num{1151794} \\
Cropland            & \num{3561}   & 205.68 & 160.88 & 500.0 & \num{732513}  \\
Built-up            & \num{1344}   & 0.00   & 0.00   & 0.0   & 0             \\
Bare / sparse       & \num{414}    & 321.57 & 424.81 & 500.0 & \num{133012}  \\
Wetland             & \num{218}    & 20.19  & 10.53  & 500.0 & \num{4402}    \\
Shrubland           & \num{9}      & 44.34  & 33.44  & 210.2 & \num{395}     \\
\bottomrule
\end{tabular}
\end{table}

\subsection{Habitat Quality: Spatial Maps and Cumulative Threats}

The full pixel-level InVEST Habitat Quality model output is mapped in
Figure~\ref{fig:hq-map}. The upper two panels render the habitat
quality index $Q$ itself on a red-to-green diverging colour ramp,
with $Q=1$ (pristine) shown in deep green and $Q=0$ (fully degraded)
in deep red. The Andaman mean Q is 0.864 and the Nicobar mean is
similar; both island groups appear overwhelmingly green-dominant,
reflecting the fact that 96\,\% of the surface is closed-canopy tree
cover or mangrove. The few red and orange pixels delineate the
settlement-around-Port-Blair cluster, the small Nicobar urban core
and the bare/grassland edges around clearings. The lower two panels
render the InVEST \emph{cumulative threat layer} $D_x$ on a sequential
red ramp; this is the proximity-weighted threat surface that
modulates the habitat-quality output via the half-saturation function
described in Section~\ref{sec:methods-habitat}. The cumulative threat
surface reveals the spatial signature of the InVEST decay kernel:
each isolated cropland or built-up pixel radiates a soft red halo
into the surrounding forest, indicating ``forest pixels close to
threats whose habitat suitability has been pulled down from the
$Q=1$ baseline''. The maximum $D_x$ on the Andamans is 0.784, on the
Nicobars somewhat lower.

\begin{figure}[H]
\centering
\includegraphics[width=0.92\textwidth]{habitat/habitat_quality_index_map.png}
\caption{InVEST habitat quality output for the Andaman \& Nicobar
Islands (2024). Top row: the habitat quality index $Q_x$ rendered on
a red-to-green diverging ramp. Bottom row: the cumulative
proximity-weighted threat layer $D_x$ (roads, built-up and cropland)
on a sequential red ramp. The halo of degraded $Q$ around each
threat source is the spatial fingerprint of the InVEST decay kernel.}
\label{fig:hq-map}
\end{figure}

\subsection{Habitat Quality by Land-Cover Class}

The pixel-level habitat quality distributions (Figure~\ref{fig:hq-violin})
follow the standard InVEST per-class ordering once they are sorted by
median: Built-up sits at $Q=0$, Cropland at $Q=0.10$, Bare/Sparse at
$Q=0.15$, Grassland at $Q=0.38$, Wetland at $Q=0.45$, Shrubland at
$Q=0.50$, Mangroves at $Q=0.90$ and Tree Cover at $Q=1.00$. Tree Cover
and Mangroves both show long downward tails towards $Q\approx 0.25$
that correspond to forest pixels close to settlement edges and
infrastructure---these are not artefacts of the violin estimator but
real signal from the InVEST threat-weighted degradation function.

\begin{figure}[H]
\centering
\includegraphics[width=0.78\textwidth]{synthesis/stat_distribution_habitat_quality.png}
\caption{Per-land-cover-class distributions of the InVEST-style
Habitat Quality Index $Q$, sorted ascending by median. Built-up sits
exactly at $Q=0$ as a modelling convention, and Shrubland's small
sample size ($n=318$) is flagged with a dashed outline.}
\label{fig:hq-violin}
\end{figure}

The bivariate relationship between Habitat Quality and Soil Erosion
(Figure~\ref{fig:hq-se-kde}) is consistent across the two largest
classes. Tree-cover pixels cluster in the high-$Q$, low-erosion
corner and cropland pixels cluster in the low-$Q$, moderate-erosion
corner, with a negative log-linear relationship across all valid
pixels ($r=-0.71$, $p<0.001$, $n=16{,}000$). The relationship is
substantively meaningful: an increase from $Q=0.10$ to $Q=1.00$
corresponds to a drop in mean soil loss of approximately a factor of
55, a quantitative confirmation of the well-documented
\citep{morgan2005} co-benefit between forest retention and soil
retention.

\begin{figure}[H]
\centering
\includegraphics[width=0.80\textwidth]{synthesis/bivariate_habitat_erosion_kde.png}
\caption{Bivariate kernel-density-estimated scatter of pixel-level
Habitat Quality versus annual soil erosion. Tree-cover pixels (green,
$n=12{,}000$) and Cropland pixels (orange, $n=4{,}000$) occupy
opposite corners; a log-linear regression yields $r=-0.71$
($p<0.001$). Marginal density curves are rendered along the top and
right edges.}
\label{fig:hq-se-kde}
\end{figure}

A horizontal bar summary of mean habitat quality per ESA WorldCover
class is shown in Figure~\ref{fig:hq-bar}. Tree cover and Mangroves
dominate the top of the ranking with $Q$ near or at the InVEST
ceiling of 1.0, while Built-up, Cropland and Bare-Sparse collapse
toward zero, exactly as expected from the lookup-table baselines.
This visual confirms the ecological ordering implied by the violin
distributions but in a form better suited to summary tables and
policy briefings.

\begin{figure}[H]
\centering
\includegraphics[width=0.80\textwidth]{habitat/habitat_quality_by_landcover.png}
\caption{Mean InVEST habitat quality $Q$ per ESA WorldCover class,
displayed as a horizontal bar chart sorted descending. Tree Cover and
Mangroves dominate the top of the ranking; Built-up, Cropland and
Bare-Sparse occupy the bottom.}
\label{fig:hq-bar}
\end{figure}

\subsection{Temporal Habitat-Quality Change (2000 $\rightarrow$ 2024)}

The change in habitat quality between the year-2000 land-cover
baseline and the year-2024 state---computed by re-running the InVEST
pipeline with the historical land-cover map and differencing the
output---is mapped in Figure~\ref{fig:hq-delta}. Two patterns are
worth highlighting. First, the change is overwhelmingly negative:
\SI{7960}{\hectare} of Andaman habitat and \SI{3715}{\hectare} of
Nicobar habitat are now in a degraded state ($\Delta Q < 0$) relative
to the year-2000 baseline, while \emph{zero} hectares register a
positive $\Delta Q$. Second, the magnitude of the $\Delta Q$
distribution in the right-hand histogram is striking: most of the
degraded pixels lose between 0.0 and $-0.2$ of quality (small
threat-halo pixels at the forest perimeter), but a heavy tail extends
all the way to $-0.9$, corresponding to forest pixels that converted
directly to built-up or bare surface and lost essentially all habitat
value. The bimodal shape of the histogram---a large peak near zero
and a smaller peak near $-0.7$---separates ``proximity degradation''
from ``conversion degradation'' as the two distinct mechanisms by
which habitat quality is being eroded.

\begin{figure}[H]
\centering
\includegraphics[width=0.96\textwidth]{habitat/habitat_quality_delta_map.png}
\caption{Temporal change in habitat quality between 2000 and 2024
($\Delta Q$). The left two maps render the per-pixel delta on a
diverging blue-red ramp; the right panel shows the distribution of
$\Delta Q$ on a log pixel-count axis. No pixels improved
(zero hectares with $\Delta Q > 0$); \SI{11675}{\hectare} are now
degraded.}
\label{fig:hq-delta}
\end{figure}

\subsection{Habitat-Quality Sensitivity Analysis}

The InVEST half-saturation distance $D_{1/2}$ and shape parameter $k$
control how aggressively threats decay with distance, and choosing
them is unavoidably judgemental. Figure~\ref{fig:hq-sens} reports a
local sensitivity analysis in which each parameter was independently
perturbed by $\pm 50\,\%$ from its default value (default $D_{1/2}=0.5$,
$k=2.5$), and the resulting mean $\overline{Q}$ was re-tabulated. The
mean is reasonably insensitive to $D_{1/2}$ (it varies by only
$\pm 0.03$ across the perturbation range) but more sensitive to $k$
(varies by $\pm 0.05$). Crucially, the qualitative ordering of
land-cover classes does not change under any parameter perturbation
within the tested range---Tree Cover is always at the top, Built-up
always at the bottom---so the synthesis conclusions about which
classes are losing habitat are robust to parameter choice. The
sensitivity envelope is a useful uncertainty bound on the
\emph{absolute} mean $Q$ but does not threaten the
\emph{relative} per-class story.

\begin{figure}[H]
\centering
\includegraphics[width=0.85\textwidth]{habitat/habitat_sensitivity_analysis.png}
\caption{Local sensitivity of the mean habitat quality $\overline{Q}$
to the InVEST half-saturation distance $D_{1/2}$ and shape parameter
$k$. Each parameter is perturbed independently by $\pm 50\,\%$ from
its default. The mean varies by less than $\pm 0.05$ across the
tested range, and the relative ranking of land-cover classes is
preserved.}
\label{fig:hq-sens}
\end{figure}

\subsection{Multifunctional Trade-offs Across Services}

Figure~\ref{fig:radar} summarises four land-cover types on three
ecosystem-service axes simultaneously. Tree cover dominates every
axis (HQ 92, Carbon 100, Soil Retention 99), Mangroves are close
behind (HQ 86, Carbon 66, Soil 97), Cropland sits in the middle on
all three axes (HQ 10, Carbon 32, Soil 56), and Built-up shows the
diagnostic ``narrow sliver'' shape that arises when habitat and
carbon collapse to zero while soil retention reads as 99 simply
because RUSLE has nothing to compute against an impervious surface.
The four radar values are computed directly from the same rasters
that produce the violin charts---no hand-tuned constants---so the
synthesis is internally consistent with the rest of the report.

\begin{figure}[H]
\centering
\includegraphics[width=0.66\textwidth]{synthesis/tradeoff_radar_chart.png}
\caption{Multidimensional ecosystem-service trade-off radar for four
representative land-cover classes. Values are normalised to the
empirical class maxima so the maximum on each spoke is exactly 100.
The Built-up sliver shape (HQ=0, Carbon=0, Soil=99) reveals the
modelling convention that RUSLE has no sediment yield to compute on
an impervious surface.}
\label{fig:radar}
\end{figure}

\subsection{Ecosystem Collapse Index and Triple-Collapse Hotspots}
\label{sec:eci-results}

The ECI map shown in Figure~\ref{fig:eci-map} renders the multifunctional
collapse signal directly. Hotspot density---visualised as hexbin counts
in Figure~\ref{fig:hexbin}---is concentrated on Middle and North
Andaman around the Saddle Peak ridge and on the southern stretch of
Great Nicobar; the spatial pattern is strikingly clustered rather than
uniformly distributed. The 95\textsuperscript{th}-percentile threshold
flags \SI{2404}{\hectare} of triple-collapse pixels (\SI{2094}{\hectare}
in the Andamans, \SI{310}{\hectare} in the Nicobars). Global Moran's
$I$ on the habitat-quality delta raster is $0.33$
($p=0.005$, 199 permutations, $n=2{,}260$ valid degradation cells)
and on the RUSLE soil-loss delta raster is $0.20$
($p=0.005$, $n=1{,}374$). The interpretation is that the spatial
configuration of degradation is significantly different from random:
hotspots are clustered rather than dispersed, which has direct
implications for prioritising management interventions geographically.

\begin{figure}[H]
\centering
\includegraphics[width=0.92\textwidth]{synthesis/eci_collapse_hotspots_map.png}
\caption{Ecosystem Collapse Index (ECI) map with triple-collapse
hotspots shown in red and the ECI value distribution shown in the
right panel. The 95\textsuperscript{th}-percentile threshold (red
dashed line in the histogram) demarcates the top-5\,\% of impacted
pixels.}
\label{fig:eci-map}
\end{figure}

\begin{figure}[H]
\centering
\includegraphics[width=0.92\textwidth]{synthesis/hotspot_hexbin_density_map.png}
\caption{Spatial hexbin clustering of triple-collapse hotspots
shown separately for the Andaman (left) and Nicobar (right) island
groups, sharing a single log-scale colour bar. The largest hotspot
density is on Middle/North Andaman near the central ridge; the
Great Nicobar pattern is offset to the south-east.}
\label{fig:hexbin}
\end{figure}

\subsection{Cross-Validation: GEDI L4B versus Saatchi 2011}
\label{sec:cv-results}

Figure~\ref{fig:cv} compares the two biomass products on the forest-
restricted pixel set (\num{2314513} pairs). The headline result is
\emph{not} a clean cross-validation: Pearson's $r$ is a modest 0.254,
Lin's CCC essentially zero (0.0213), RMSE 168.8 Mg/ha and bias
$+152.6$ Mg/ha---meaning that on average GEDI reports about 153 Mg/ha
more biomass than Saatchi reports for the same forest pixel. The
mean ratio is 3.45$\times$ and the RMA slope is 4.56. Log--log Pearson
$r$ rises to 0.36, showing that the two products rank pixels in a
similar relative order even though their absolute scales diverge.

\begin{figure}[H]
\centering
\includegraphics[width=0.96\textwidth]{validation/agb_cross_validation_upgraded.png}
\caption{Inter-product comparison of GEDI L4B versus Saatchi 2011
biomass on the forest-restricted pixel set ($n=2{,}314{,}513$).
The hexbin scatter (left) shows the cloud sitting well above the 1:1
line, the overlaid histograms (centre) show GEDI's median of 216 Mg/ha
versus Saatchi's median of 62 Mg/ha, and the metrics table (right)
reports bootstrap 95\,\% confidence intervals on every quantity. The
inset note clarifies that this is an inter-product comparison rather
than a ground-truth validation.}
\label{fig:cv}
\end{figure}

A diagnostic test that rescaled Saatchi by a factor of 2.13 (the
literature conversion from Mg C/ha to Mg AGB/ha) closed only half of
the bias gap, indicating that the discrepancy is not a simple
unit-conversion error. The remaining factor reflects two well-known
features of the Saatchi 2011 product: it was calibrated against
GLAS-era 2003--2007 GLAS waveforms, and it was conservative on tall
closed-canopy tropical forest \citep{mitchard2014}. The
inter-comparison should therefore be read as evidence of
product-level divergence, not as a validation of either product
against ground truth.

\subsection{Supplementary Ecosystem Services}

The supplementary services module (Figure~\ref{fig:supp}) extends the
biophysical accounting in four directions.
Coastal protection: \SI{54097}{\hectare} of protective ($CPI>0.5$)
mangrove area corresponds to an annual avoided-damage value of
\$121.2 million using the 10\,\%-of-global Menendez willingness-to-pay
benchmark; the 15-year discounted NPV of high-value mangrove lost
(\SI{971}{\hectare}, $CPI>0.5$) is \$24.2 million.
Freshwater yield: the Budyko-style P--AET partition gives
\SI{3140}{\giga\litre\per\year} of water yield across the displayed
classes, dominated by tree cover (2681 GL); counterfactual
forest-to-cropland conversion would increase runoff by \SI{3574}{\giga\litre\per\year}
but with the well-known side effect of higher flood and erosion risk.
Pollination: the Lonsdorf $S$ index ranges from 0.76 on tree cover to
0.00 on built-up, with cropland at 0.20; the per-hectare loss going
from tree cover to cropland is $\Delta S = -0.565$ and
\SI{11016}{unit\hectare} of suitability has effectively been lost
over the deforested area to date.
Soil organic carbon: the hybrid SoilGrids-plus-class-mean baseline
gives a mean SOC of \SI{75.3}{\mega\gram\of\carbon\per\hectare}; a 25\,\% loss factor
\citep{don2011} applied to the deforested pixel set gives
\SI{365}{\kilo\tonne\of\carbon} lost, or about \SI{1341}{\giga\gram\of\carbondioxide} of CO$_2$e
emissions on top of the aboveground budget---roughly +96\,\% of the
main carbon-loss figure---valued at \$68.4 million at the social cost
of carbon.

\begin{figure}[H]
\centering
\includegraphics[width=0.95\textwidth]{supplementary/supplementary_services_panel.png}
\caption{Supplementary ecosystem-service module summarising four
additional services not part of the main carbon-soil-habitat budget:
coastal protection (mangrove buffer index, upper-left), freshwater
yield by land-cover class (upper-right), Lonsdorf-style pollination
suitability (lower-left) and soil-organic-carbon loss from
deforestation (lower-right). Together these raise the projected 2040
economic damage from a lower-bound \$502 M to an upper-bound
\$594 M.}
\label{fig:supp}
\end{figure}

\subsection{Predictive Forecasting: Forest Cover to 2040}
\label{sec:forecast-forest}

Extrapolating the 2001--2023 forest-loss trajectory forward gives a
projected additional \SI{25056}{\hectare} of deforestation by 2040
under the business-as-usual scenario, calculated as the historical
mean annual loss (\SI{812.6}{\hectare\per\year}) multiplied by the
17-year forecast horizon and modulated by the Sen-slope adjustment.
Figure~\ref{fig:forecast-forest} maps the projected loss footprint
spatially: the new loss pixels (red) are sited preferentially along
the existing road network and at the perimeter of current cropland
clusters, because the underlying probabilistic land-use model
parameterises loss probability as a function of proximity to existing
disturbance. The total forecast loss represents an additional 12\,\%
of the standing forest area in the affected zones and would push the
cumulative carbon-equivalent debt past 2.3 Gg CO$_2$e by 2040 under
BAU assumptions.

\begin{figure}[H]
\centering
\includegraphics[width=0.92\textwidth]{predictive/forecast_forest_cover_2040.png}
\caption{Projected forest-cover loss footprint for ANI by 2040 under
the business-as-usual scenario. Historical 2001--2024 loss is shown
in orange, projected new 2024--2040 loss in red. The new-loss
pattern follows the existing road and settlement-edge geography
because the underlying model treats proximity to existing disturbance
as the dominant explanatory variable.}
\label{fig:forecast-forest}
\end{figure}

\subsection{Three-Scenario Predictive Forest-Collapse Map}
\label{sec:forecast-tri}

To bracket the deep uncertainty in future deforestation pressure,
three counterfactual scenarios were run forward to 2060: a
\emph{Conservation} scenario in which the historical loss rate drops
to $0.2\times$ its 2001--2024 mean (consistent with full enforcement
of core-protected status), a \emph{Business-As-Usual} scenario at
$1.0\times$, and an \emph{Escalation} scenario at $2.5\times$
(consistent with regional infrastructure projects of the type
announced for Great Nicobar). Figure~\ref{fig:forecast-tri} renders
the three predicted forest-cover endings in side-by-side panels.
The Conservation scenario adds just \SI{6630}{\hectare} of new loss
and costs $\sim$\$248 M in cumulative damages, BAU adds
\SI{52580}{\hectare} ($\sim$\$1{,}964 M), and Escalation adds
\SI{80015}{\hectare} ($\sim$\$2{,}988 M)---roughly four-fold and
twelve-fold the BAU loss respectively. The Nicobar group is
proportionally far more vulnerable than the Andaman group under
Escalation: Great Nicobar's forested fraction drops from $\sim$80\,\%
to under 40\,\% in the Escalation panel, while the equivalent
Andaman drop is only $\sim$10 percentage points.

\begin{figure}[H]
\centering
\includegraphics[width=0.96\textwidth]{predictive/forecast_tri_scenario_2060.png}
\caption{Three-scenario predictive forecast 2024--2060.
(A) Conservation (0.2$\times$ historic rate, core-protected),
(B) Business-As-Usual (1.0$\times$ historic rate), and
(C) Escalation (2.5$\times$ historic rate). Each panel shows the
surviving forest backdrop (faded green), the historical 2000--2024
loss footprint (yellow) and the new projected 2024--2060 loss (red).
Totals at the foot of each panel report the additional new loss and
the cumulative damage cost.}
\label{fig:forecast-tri}
\end{figure}

\subsection{Multi-Scenario Projected Economic Damages to 2060}

The cumulative economic-damage trajectories for the three scenarios
are plotted side-by-side through 2024--2060 in
Figure~\ref{fig:forecast-econ}. Conservation asymptotes at
\$248 M---the trajectory levels off because once the core-protected
zone is enforced, almost no new conversion occurs and only the
already-incurred damage accumulates. BAU continues approximately
linearly to \$1,964 M by 2060. Escalation rises steeply to
\$2,988 M by 2060 and shows no sign of saturating. The
``avoided-damage band'' annotations on the figure quantify the
implicit value of switching between scenarios: moving from
Escalation to BAU saves \$1,024 M; moving from BAU to Conservation
saves a further \$1,716 M; the total avoided damage of going from
Escalation to Conservation is \$2,741 M---i.e. the maximum
willingness-to-pay for a Conservation policy in net-present-value
terms, against an Escalation counterfactual, is approximately
\$2.7 billion over the 2024--2060 horizon.

\begin{figure}[H]
\centering
\includegraphics[width=0.90\textwidth]{predictive/forecast_economic_damages_2060.png}
\caption{Cumulative projected economic damages 2024--2060 under the
three scenarios. Conservation (green) saturates around \$248 M, BAU
(orange) continues linearly to \$1,964 M, Escalation (red) reaches
\$2,988 M. Avoided-damage call-outs in the middle of the figure
quantify the value of moving between scenarios.}
\label{fig:forecast-econ}
\end{figure}

\subsection{Economic Damage Projection to 2040}

Combining the historic deforestation velocity of
\SI{812.6}{\hectare\per\year} with the three primary damage categories
gives a projected 2040 business-as-usual economic damage of
approximately \$502 million broken down as
\$101 M carbon damage (at \$51/tCO$_2$),
\$100 M dredging cost (sediment loss at \$5/tonne) and
\$301 M habitat-replacement cost (\$3000/ha for the projected
\SI{25056}{\hectare} of additional deforestation by 2040). Adding the
supplementary coastal-protection NPV (\$24.2 M) and the SOC damage
(\$68.4 M) raises the upper-bound total to \$594 million, an 18\,\%
uplift over the lower bound. The 2040 figure is the natural
midpoint of the 2060 trajectories displayed in
Figure~\ref{fig:forecast-econ}; the underlying scenarios diverge
sharply by then, with the Conservation case asymptotic at
$\sim$\$500 M, BAU at around \$1.0 billion and Escalation
approaching \$2.0 billion. These are first-order figures intended to
bracket the order of magnitude of foregone value rather than precise
predictions, and they are dominated by the highly uncertain
habitat-replacement cost component.

% ---- TABLE: HEADLINE NUMBERS ----------------------------------------
\begin{table}[H]
\centering
\caption{Headline biophysical and economic results summarised across
the 2001--2024 study period.}
\label{tab:headline}
\small
\begin{tabular}{p{6.5cm} l}
\toprule
\rowcolor{forestgreen!15}
\textbf{Quantity} & \textbf{Value (95\,\% CI where given)} \\
\midrule
Total forest loss (2001--2023)        & \SI{19502}{\hectare} (12351--27826) \\
Total CO$_2$e loss                    & \SI{1390}{\giga\gram} (936--1887) \\
Mann--Kendall trend $Z$ (CO$_2$e)     & $-2.51$ ($p=0.012$, decreasing) \\
Sen slope (CO$_2$e)                   & $-2.67$ Gg/yr ($-5.25$ to $-0.66$) \\
Mann--Kendall trend $Z$ (area)        & $-3.35$ ($p=0.0008$, decreasing) \\
Sen slope (area)                      & $-54.1$ ha/yr ($-83.8$ to $-20.4$) \\
\addlinespace
RUSLE median, Tree Cover              & \SI{2.83}{\tonne\per\hectare\per\year} \\
RUSLE median, Bare/Sparse             & \SI{424.81}{\tonne\per\hectare\per\year} \\
RUSLE median, Cropland                & \SI{160.88}{\tonne\per\hectare\per\year} \\
\addlinespace
Habitat Quality (Tree Cover)          & median 1.00 \\
Habitat Quality (Mangroves)           & median 0.90 \\
Habitat Quality (Cropland)            & median 0.10 \\
\addlinespace
ECI critical hotspot area (Andamans)  & \SI{2094}{\hectare} \\
ECI critical hotspot area (Nicobars)  & \SI{310}{\hectare} \\
Moran's $I$ (habitat $\Delta$)        & 0.33 ($p=0.005$, $n=2260$) \\
Moran's $I$ (RUSLE $\Delta$)          & 0.20 ($p=0.005$, $n=1374$) \\
\addlinespace
GEDI vs Saatchi Pearson $r$           & 0.254 (0.250--0.259) \\
GEDI vs Saatchi Lin CCC               & 0.0213 (0.021--0.022) \\
GEDI vs Saatchi RMSE                  & 168.76 Mg/ha (168.45--169.11) \\
GEDI vs Saatchi bias                  & $+152.59$ Mg/ha \\
\addlinespace
Economic damage 2040 (lower bound)    & \$501.8 M \\
Economic damage 2040 (upper bound)    & \$594.4 M ($+18\,\%$) \\
\bottomrule
\end{tabular}
\end{table}

% =====================================================================
% 4. ACKNOWLEDGEMENTS
% =====================================================================
\section{Acknowledgements}

This work uses freely-available remote-sensing data and would not have
been possible without the institutions and missions listed in
Section~\ref{sec:data}. In particular the ESA WorldCover Consortium,
NASA's Global Ecosystem Dynamics Investigation (GEDI) team, the
University of Maryland and World Resources Institute Global Forest
Watch partnership, the UCSB Climate Hazards Center (CHIRPS), the
ISRIC SoilGrids project and the USGS/NASA SRTM mission are all
gratefully acknowledged for the open data they produce. The Saatchi
pan-tropical biomass map distributed under the NASA Earthdata licence
was used as an independent cross-validation reference. The Survey of
India administrative boundary shapefile was used to delimit the study
area. Python, NumPy, SciPy, matplotlib, pandas, rasterio and
GeoPandas underpinned every analytical step, and I thank their
respective open-source communities. Any errors of interpretation are
mine alone.

% =====================================================================
% 5. REFERENCES
% =====================================================================
\section{References}
\renewcommand{\refname}{}      % hide thebibliography's auto-heading
\vspace{-2.2em}
\begin{thebibliography}{99}
\setlength{\itemsep}{2pt}

\bibitem[Anselin(1995)]{anselin1995}
Anselin, L. (1995).
Local Indicators of Spatial Association---LISA.
\textit{Geographical Analysis}, 27(2), 93--115.

\bibitem[Borrelli et al.(2020)]{borrelli2020}
Borrelli, P., Robinson, D.A., Panagos, P., et al. (2020).
Land use and climate change impacts on global soil erosion by water
(2015--2070).
\textit{Proceedings of the National Academy of Sciences}, 117(36),
21994--22001.

\bibitem[Don et al.(2011)]{don2011}
Don, A., Schumacher, J., \& Freibauer, A. (2011).
Impact of tropical land-use change on soil organic carbon stocks---a
meta-analysis.
\textit{Global Change Biology}, 17(4), 1658--1670.

\bibitem[Donato et al.(2011)]{donato2011}
Donato, D.C., Kauffman, J.B., Murdiyarso, D., Kurnianto, S.,
Stidham, M., \& Kanninen, M. (2011).
Mangroves among the most carbon-rich forests in the tropics.
\textit{Nature Geoscience}, 4, 293--297.

\bibitem[Dubayah et al.(2020)]{dubayah2020}
Dubayah, R., Blair, J.B., Goetz, S., et al. (2020).
The Global Ecosystem Dynamics Investigation: high-resolution laser
ranging of the Earth's forests and topography.
\textit{Science of Remote Sensing}, 1, 100002.

\bibitem[Dubayah et al.(2022)]{dubayah2022}
Dubayah, R., Armston, J., Kellner, J.R., et al. (2022).
GEDI L4A Footprint Level Aboveground Biomass Density, Version 2.1.
ORNL DAAC, Oak Ridge, Tennessee, USA.

\bibitem[Efron \& Tibshirani(1993)]{efron1993}
Efron, B. \& Tibshirani, R.J. (1993).
\textit{An Introduction to the Bootstrap}.
Chapman \& Hall/CRC.

\bibitem[Egoh et al.(2008)]{egoh2008}
Egoh, B., Reyers, B., Rouget, M., et al. (2008).
Mapping ecosystem services for planning and management.
\textit{Agriculture, Ecosystems \& Environment}, 127(1--2), 135--140.

\bibitem[EPA(2023)]{epa2023}
United States Environmental Protection Agency (2023).
EPA Report on the Social Cost of Greenhouse Gases: Estimates
Incorporating Recent Scientific Advances. EPA, Washington DC.

\bibitem[Hansen et al.(2013)]{hansen2013}
Hansen, M.C., Potapov, P.V., Moore, R., et al. (2013).
High-Resolution Global Maps of 21st-Century Forest Cover Change.
\textit{Science}, 342, 850--853.

\bibitem[Harris et al.(2021)]{harris2021}
Harris, N.L., Gibbs, D.A., Baccini, A., et al. (2021).
Global maps of twenty-first century forest carbon fluxes.
\textit{Nature Climate Change}, 11, 234--240.

\bibitem[Kendall(1975)]{kendall1975}
Kendall, M.G. (1975).
\textit{Rank Correlation Methods}. Griffin, London.

\bibitem[Koh et al.(2016)]{koh2016}
Koh, I., Lonsdorf, E.V., Williams, N.M., et al. (2016).
Modeling the status, trends, and impacts of wild bee abundance in the
United States.
\textit{Proceedings of the National Academy of Sciences}, 113(1),
140--145.

\bibitem[Lin(1989)]{lin1989}
Lin, L.I. (1989).
A concordance correlation coefficient to evaluate reproducibility.
\textit{Biometrics}, 45(1), 255--268.

\bibitem[Mann(1945)]{mann1945}
Mann, H.B. (1945).
Nonparametric tests against trend.
\textit{Econometrica}, 13(3), 245--259.

\bibitem[Menendez et al.(2020)]{menendez2020}
Menendez, P., Losada, I.J., Torres-Ortega, S., Narayan, S., \&
Beck, M.W. (2020).
The global flood-protection benefits of mangroves.
\textit{Scientific Reports}, 10, 4404.

\bibitem[Mitchard et al.(2014)]{mitchard2014}
Mitchard, E.T.A., Saatchi, S.S., Baccini, A., et al. (2014).
Markedly divergent estimates of Amazon forest carbon density from
ground plots and satellites.
\textit{Global Ecology and Biogeography}, 23, 935--946.

\bibitem[Montgomery(2007)]{montgomery2007}
Montgomery, D.R. (2007).
Soil erosion and agricultural sustainability.
\textit{Proceedings of the National Academy of Sciences}, 104(33),
13268--13272.

\bibitem[Moore et al.(1992)]{moore1992}
Moore, I.D., Grayson, R.B., \& Ladson, A.R. (1991).
Digital terrain modelling: a review of hydrological, geomorphological
and biological applications.
\textit{Hydrological Processes}, 5, 3--30.

\bibitem[Moran(1950)]{moran1950}
Moran, P.A.P. (1950).
Notes on continuous stochastic phenomena.
\textit{Biometrika}, 37(1--2), 17--23.

\bibitem[Morgan(2005)]{morgan2005}
Morgan, R.P.C. (2005).
\textit{Soil Erosion and Conservation}, 3rd ed.
Blackwell Publishing.

\bibitem[Nelson et al.(2009)]{nelson2009}
Nelson, E., Mendoza, G., Regetz, J., et al. (2009).
Modeling multiple ecosystem services, biodiversity conservation,
commodity production, and tradeoffs at landscape scales.
\textit{Frontiers in Ecology and the Environment}, 7(1), 4--11.

\bibitem[Panagos et al.(2015)]{panagos2015}
Panagos, P., Borrelli, P., Poesen, J., et al. (2015).
The new assessment of soil loss by water erosion in Europe.
\textit{Environmental Science \& Policy}, 54, 438--447.

\bibitem[Powers et al.(2018)]{powers2018}
Powers, J.S., Corre, M.D., Twine, T.E., \& Veldkamp, E. (2018).
Geographic bias of field observations of soil carbon stocks with
tropical land-use changes precludes spatial extrapolation.
\textit{Proceedings of the National Academy of Sciences}, 108(15),
6318--6322.

\bibitem[Rao et al.(2018)]{rao2018}
Rao, N.D., Min, J., DeFries, R., et al. (2018).
Ecosystem-service bundles and their spatial determinants.
\textit{Ecosystem Services}, 30, 102--113.

\bibitem[Renard et al.(1997)]{renard1997}
Renard, K.G., Foster, G.R., Weesies, G.A., et al. (1997).
\textit{Predicting Soil Erosion by Water: A Guide to Conservation
Planning with the Revised Universal Soil Loss Equation (RUSLE)}.
USDA Agriculture Handbook 703.

\bibitem[Saatchi et al.(2011)]{saatchi2011}
Saatchi, S.S., Harris, N.L., Brown, S., et al. (2011).
Benchmark map of forest carbon stocks in tropical regions across
three continents.
\textit{Proceedings of the National Academy of Sciences}, 108(24),
9899--9904.

\bibitem[Sen(1968)]{sen1968}
Sen, P.K. (1968).
Estimates of the regression coefficient based on Kendall's tau.
\textit{Journal of the American Statistical Association}, 63,
1379--1389.

\bibitem[Sharp et al.(2020)]{sharp2020}
Sharp, R., Tallis, H.T., Ricketts, T., et al. (2020).
\textit{InVEST 3.9.0 User's Guide}.
The Natural Capital Project, Stanford University.

\bibitem[Sun et al.(2006)]{sun2006}
Sun, G., Zhou, G., Zhang, Z., et al. (2006).
Potential water yield reduction due to forestation across China.
\textit{Journal of Hydrology}, 328(3--4), 548--558.

\bibitem[Terrado et al.(2016)]{terrado2016}
Terrado, M., Sabater, S., Chaplin-Kramer, B., et al. (2016).
Model development for the assessment of terrestrial and aquatic
habitat quality in conservation planning.
\textit{Science of the Total Environment}, 540, 63--70.

\bibitem[Theil(1950)]{theil1950}
Theil, H. (1950).
A rank-invariant method of linear and polynomial regression
analysis.
\textit{Proceedings of the Royal Netherlands Academy of Sciences},
53, 386--392.

\bibitem[Warton et al.(2006)]{warton2006}
Warton, D.I., Wright, I.J., Falster, D.S., \& Westoby, M. (2006).
Bivariate line-fitting methods for allometry.
\textit{Biological Reviews}, 81(2), 259--291.

\bibitem[Westlake(1963)]{westlake1963}
Westlake, D.F. (1963).
Comparisons of plant productivity.
\textit{Biological Reviews}, 38, 385--425.

\bibitem[Wischmeier \& Smith(1978)]{wischmeier1978}
Wischmeier, W.H. \& Smith, D.D. (1978).
\textit{Predicting Rainfall Erosion Losses---A Guide to Conservation
Planning}. USDA Agriculture Handbook 537.

\bibitem[Zhang et al.(2001)]{zhang2001}
Zhang, L., Dawes, W.R., \& Walker, G.R. (2001).
Response of mean annual evapotranspiration to vegetation changes at
catchment scale.
\textit{Water Resources Research}, 37(3), 701--708.

\end{thebibliography}

% =====================================================================
% APPENDIX
% =====================================================================
\appendix

\section{Code and Data Layout}

The project is organised under a single working directory with three
main top-level folders: \texttt{data/} (raw and processed
GeoTIFFs), \texttt{results/} (CSV outputs and intermediate rasters)
and \texttt{figures/} (publication-ready PNGs grouped by analytical
theme). The analysis pipeline is implemented in Python 3.11 with
NumPy, SciPy, pandas, matplotlib, rasterio and GeoPandas. The
ordered execution sequence is: (1) \texttt{src/carbon\_analysis.py}
to build the GEDI biomass baseline and the annual carbon-loss CSV;
(2) \texttt{src/soil\_retention.py} to compute the four RUSLE factor
rasters, the annual soil-loss raster and the per-class summary CSV;
(3) \texttt{src/habitat\_quality.py} for the InVEST-style
habitat-quality index and its per-class summary; (4)
\texttt{src/synthesis\_hotspots.py} for the ECI raster and the
top-5\,\% hotspot mask; (5)
\texttt{src/supplementary\_services.py} for the four supplementary
modules and the updated economic upper bound; (6)
\texttt{src/validation\_stats.py} for the Mann--Kendall, Theil--Sen,
bootstrap cross-validation and Moran's $I$ tests; and (7) the
\texttt{src/render\_synthesis\_light.py} and other rendering scripts
to produce the publication figures embedded in this report.

\section{Reproducibility}

All scripts are deterministic given the same input rasters and the
NumPy random seed 42 used throughout. The validation summary JSON
(\texttt{results/validation\_summary.json}) and the flat-CSV summary
(\texttt{results/validation\_metrics.csv}) record every numeric
quantity reported in this document.

\end{document}
